% =====================================================================================
% FICHE DE CALCULS -- tout ce que le memoire CONSERVE, avec la formule, les entrees et la
% sortie. 13 aout 2026.
%
% OBJET. Les deux notes de suivi disent ce qui est integre et ce qui ne l'est pas. Celle-ci
% dit COMMENT c'est calcule : une formule, ses entrees, sa sortie, et le script qui l'imprime.
% Elle est faite pour etre ouverte pendant un point quand une question porte sur un chiffre.
%
% REGLE TENUE PARTOUT : aucun nombre de cette fiche n'est calcule a la main. Chacun est
% imprime par un script dont la sortie est versionnee, et le harnais le verifie. Un ratio
% calcule pendant la redaction n'est verifiable par personne.
%
% CHAQUE SECTION CITE SES SCRIPTS, faute de quoi le harnais la compte hors controle : c'est
% la faille qui avait laisse passer un intervalle faux sur une slide de juillet.
% =====================================================================================
\documentclass[11pt,a4paper]{article}

\usepackage[utf8]{inputenc}
\usepackage[T1]{fontenc}
\usepackage{lmodern}
\usepackage[french]{babel}
\usepackage[margin=1.9cm]{geometry}
\usepackage{microtype}
\usepackage{amsmath, amssymb}
\usepackage{array}
\usepackage{booktabs}
\usepackage{longtable}
\usepackage[table]{xcolor}
\usepackage[most]{tcolorbox}
\usepackage{fancyhdr}
\usepackage{lastpage}
\usepackage{eurosym}

\definecolor{navy}{HTML}{1F3864}
\definecolor{lightblue}{HTML}{D6E0F0}
\definecolor{softgray}{HTML}{F4F4F2}

\newtcolorbox{cle}[1][]{colback=lightblue!40,colframe=navy,boxrule=0.6pt,
  arc=2pt,left=8pt,right=8pt,top=6pt,bottom=6pt,#1}
\newtcolorbox{res}[1][]{colback=softgray,colframe=black!30,boxrule=0.5pt,
  arc=2pt,left=7pt,right=7pt,top=5pt,bottom=5pt,#1}

\newcommand{\src}[1]{\hfill{\footnotesize\textsf{script~\texttt{#1}}}}

\pagestyle{fancy}
\fancyhf{}
\lfoot{\footnotesize Fiche de calculs --- 13 août 2026}
\rfoot{\footnotesize page \thepage/\pageref{LastPage}}
\renewcommand{\headrulewidth}{0pt}

\title{\vspace{-1.3cm}\color{navy}\bfseries Fiche de calculs\\[2pt]
  \large Tout ce que le mémoire conserve : la formule, les entrées, la sortie}
\author{Kélian Kaddouri}
\date{13 août 2026}

\begin{document}
\maketitle
\thispagestyle{fancy}
\vspace{-0.7cm}

% HARNAIS-HORS-SECTION: le preambule est un mode d'emploi, pas un resultat. Les quelques
% nombres qu'il contient decrivent le dispositif de verification lui-meme (combien de chiffres
% perimes ont ete trouves de telle facon) et ne sont produits par aucun script du pipeline.
\begin{cle}
\textbf{Comment se servir de cette fiche.} Un bloc par résultat conservé, dans l'ordre de la
chaîne de calcul : la sévérité, la fréquence, la cascade, l'écart entre états, le retour sur
conformité, puis les mesures de robustesse. Chaque bloc donne la \textbf{formule}, ses
\textbf{entrées} et sa \textbf{sortie}, avec le script qui l'imprime en marge.

\textbf{Une règle est tenue partout et elle est le point de méthode de ce document} : aucun
nombre n'est calculé à la main. Chacun est imprimé par un script dont la sortie est versionnée,
et un harnais vérifie que le nombre publié se retrouve bien dans la sortie du script que la
section cite. Un ratio calculé pendant la rédaction n'est vérifiable par personne, et plusieurs
des chiffres périmés trouvés cette année étaient exactement de cette nature.
\end{cle}

% =====================================================================================
\section{La sévérité d'un sinistre}
\src{\texttt{07}, \texttt{47}, \texttt{63}}

\paragraph{Le modèle.} Au-delà d'un seuil $u$, les dépassements suivent une loi de Pareto
généralisée. Le quantile de sévérité d'un \emph{sinistre} s'écrit en forme fermée
\begin{equation}
\mathrm{q}_{\alpha} \;=\; u \;+\; \frac{\sigma}{\xi}
  \left[\left(\frac{1-\alpha}{p_u}\right)^{-\xi} - 1\right],
\end{equation}
où $p_u$ est le taux de dépassement du seuil.

\begin{res}
\textbf{Entrées (calibration gelée)} : $u = 20{,}03$~M\euro, $\xi = 0{,}5954$,
$\sigma = 57{,}97$, $p_u = 0{,}1509$, sur $91$ excès.\\
\textbf{Sortie} : $\mathrm{q}_{99,5\%} = 662{,}78$~M\euro, d'intervalle à $90\,\%$
$[411{,}5\,;1\,037{,}1]$, soit un \textbf{facteur $2{,}5$} entre les deux bornes.\\
\textbf{Adéquation} : Kolmogorov-Smirnov $D = 0{,}043$, $p = 0{,}889$ ; non rejeté à $5\,\%$.
\end{res}

\paragraph{Deux précautions qui accompagnent ce nombre.} La première est de notation :
$662{,}78$ est un quantile \emph{par sinistre} et non un capital. Le confondre avec le capital
agrégé fait lire une différence d'échelle comme une contradiction, ce qui est arrivé en séance.
La seconde est que l'intervalle de $\xi$ vaut $[0{,}304\,;0{,}831]$ par bootstrap et
$[0{,}320\,;0{,}871]$ par delta-méthode : \textbf{ce sont deux intervalles distincts}, et le
mémoire publie le premier. Toute citation mélangeant une borne de l'un avec une borne de l'autre
est fautive.

% =====================================================================================
\section{La fréquence des sinistres}
\src{\texttt{08b}, \texttt{41}}

\paragraph{Le modèle.} Le nombre annuel de sinistres suit une binomiale négative de moyenne
$\lambda$ et d'indice de dispersion $\varphi$, paramétrée par
\begin{equation}
N \sim \mathcal{BN}(r,\,p), \qquad r = \frac{\lambda}{\varphi - 1},
\qquad p = \frac{r}{r + \lambda}.
\end{equation}

\begin{res}
\textbf{Entrées} : $\lambda_{\text{ref}} = 21{,}6$ sinistres par an au secteur,
$\varphi = 9{,}20$ posé.\\
\textbf{Mesures d'échelle} : la fréquence observée vaut $0{,}0988$ par firme et par an sur
l'ensemble du panel, et $0{,}2102$ sur les firmes à au moins dix événements, profil plus proche
d'une entité soumise au règlement.
\end{res}

\paragraph{L'indice de dispersion n'est pas invariant d'échelle}, et c'est ce qui réconcilie
trois chiffres qui paraissent se contredire : $\varphi = 1 + \lambda/r$ vaut $1{,}19$ sur les
cellules firme-année, $2{,}24$ sur la série annuelle détendancée, et $9{,}20$ au moteur où il est
un choix de prudence. Une échelle de lecture, une agrégation, une marge.

% =====================================================================================
\section{La cascade entre piliers}
\src{\texttt{03}, \texttt{74}}

\paragraph{La matrice de propagation.} Elle se déduit de la table de transitions par une
normalisation de Leontief
\begin{equation}
W(g) \;=\; g \cdot \frac{\mathrm{TRANS}}{\max_j s_j},
\qquad s_j = \sum_k \mathrm{TRANS}_{jk},
\end{equation}
où $s_j$ est l'\textbf{émission} du pilier $j$, non sa réception. La convention est
$W_{jk}$ de la source $j$ vers la cible $k$ : \emph{la ligne émet, la colonne reçoit}.

\begin{res}
\textbf{Pourquoi le diviseur est l'émission.} La normalisation existe pour garantir que le
pilier le plus prolifique engendre \emph{au plus} $g$ descendants directs. Avec le diviseur de
réception ($2{,}3$) on obtenait $0{,}90 \times 2{,}60 / 2{,}3 = 1{,}02$, donc plus que
$g = 0{,}90$ : la borne que la normalisation est censée poser était franchie. Avec le diviseur
d'émission ($2{,}60$) on a exactement $0{,}90$.\\
\textbf{Sorties} : $\rho(W) = 0{,}506 = 0{,}562 \times g$, sous-critique ;
$R_0 = 0{,}054$ ; $\rho(W) = 1$ à $g = 1{,}78$, hors domaine.
\end{res}

\paragraph{La loi du nombre de piliers touchés est exacte.} Elle s'obtient par énumération de la
progéniture du noyau auto-évitant, sans simulation, en pondérant les amorces par leur propension.

\begin{center}\footnotesize
\renewcommand{\arraystretch}{1.15}
\begin{tabular}{@{}c c c@{}}
\toprule
\textbf{Piliers touchés} & \textbf{Conforme} ($g=0{,}45$) & \textbf{Non conforme} ($g=0{,}90$) \\
\midrule
$1$ & $68{,}85\,\%$ & $37{,}69\,\%$ \\
$2$ & $25{,}10\,\%$ & $38{,}10\,\%$ \\
$3$ & $5{,}31\,\%$ & $18{,}29\,\%$ \\
$4$ & $0{,}70\,\%$ & $5{,}22\,\%$ \\
$5$ & $0{,}04\,\%$ & $0{,}69\,\%$ \\
\midrule
\emph{plus d'un pilier} & $\mathbf{31{,}15\,\%}$ & $\mathbf{62{,}31\,\%}$ \\
\emph{moyenne} & $1{,}380$ & $1{,}931$ \\
\bottomrule
\end{tabular}
\end{center}

% =====================================================================================
\section{L'écart de capital entre états de conformité}
\src{\texttt{43}, \texttt{68}, \texttt{74}}

\paragraph{Les quatre canaux, et ils ne s'additionnent pas.} La non-conformité n'ajoute aucune
pénalité au capital : elle déplace quatre paramètres de la loi de perte.

\begin{center}\footnotesize
\renewcommand{\arraystretch}{1.15}
\begin{tabular}{@{}l c c l@{}}
\toprule
\textbf{Canal} & \textbf{Conforme} & \textbf{Non conforme} & \textbf{Statut} \\
\midrule
Fréquence $\lambda$ & $21{,}6$ & $53{,}6$ & calibrable \\
Détection $p_u$ & $0{,}128$ & $0{,}181$ & calibrable \\
Propagation $g$ & $0{,}45$ & $0{,}90$ & borné \\
Accumulation $\varphi_{cs}$ & $0$ & $0{,}68$ & borné \\
\bottomrule
\end{tabular}
\end{center}

\begin{res}
\textbf{Sortie} : capital de $6\,049$ à $20\,188$~M\euro, soit un écart
$\Delta_{\text{DORA}} = 14\,139$~M\euro{} et un \textbf{facteur $3{,}34$} entre les deux états.
Bruit de simulation de l'écart, quatre graines appariées : $\pm 734$~M\euro.\\
\textbf{Résolution} : $40\,000$ années par graine, à sévérité de base et échelle inchangées ;
seuls les quatre canaux bougent.
\end{res}

\paragraph{Où passent les cinq milliards manquants.} Sommer les quatre canaux pris isolément
donne $9\,138$, non $14\,139$. L'écart n'est pas une erreur, c'est l'\textbf{interaction}. On
l'établit en énumérant les $2^4$ configurations du treillis et en prenant la transformée de
Möbius
\begin{equation}
m(S) \;=\; \sum_{T \subseteq S} (-1)^{|S| - |T|}\, v(T),
\qquad \sum_{S \neq \emptyset} m(S) \;=\; v(\{1,2,3,4\}).
\end{equation}

\begin{res}
\textbf{Sortie par ordre} : ordre~1 $9\,138$, ordre~2 $5\,345$, ordre~3 $-691$, ordre~4 $+346$.
La somme des quinze termes redonne $14\,139$ \textbf{à la précision machine}.\\
\textbf{Interaction} : $14\,139 - 9\,138 = 5\,001$~M\euro, soit $+35\,\%$.
\end{res}

\paragraph{Trois lectures d'un même canal, et elles ne sont pas interchangeables.} Isolée depuis
l'état conforme, fermeture depuis l'état non conforme, et Shapley.

\begin{center}\footnotesize
\renewcommand{\arraystretch}{1.15}
\begin{tabular}{@{}l r r r@{}}
\toprule
\textbf{Canal} & \textbf{Isolé} & \textbf{Fermeture} & \textbf{Shapley} \\
\midrule
Fréquence & $4\,328 \pm 463$ & $\mathbf{8\,775} \pm 1\,004$ & $6\,546 \pm 457$ \\
Détection & $1\,448 \pm 277$ & $4\,562 \pm 819$ & $2\,928 \pm 343$ \\
Accumulation P4 & $1\,728 \pm 210$ & $3\,402 \pm 493$ & $2\,576 \pm 222$ \\
Propagation $W$ & $1\,633 \pm 188$ & $2\,402 \pm 442$ & $2\,088 \pm 152$ \\
\midrule
\textbf{somme} & $9\,138$ & $19\,141$ & $\mathbf{14\,139}$ \\
\bottomrule
\end{tabular}
\end{center}

\vspace{2pt}
\noindent La colonne isolée \emph{manque} $5\,001$, celle de fermeture le \emph{dépasse}
d'autant. C'est la colonne de \textbf{fermeture} qu'un plan de remédiation doit citer : remédier
un canal rapporte davantage à une entité défaillante partout, puisque le canal ferme aussi les
croisés qu'il portait. Shapley est la seule colonne additive, mais c'est une \emph{convention}
d'attribution, et deux canaux sur quatre étant bornés, leur part n'est pas un budget.

\paragraph{Les six effets croisés, avec leur bruit.} freq$\times$det $1\,739 \pm 708$,
freq$\times$accum $1\,666 \pm 354$, freq$\times$prop $1\,182 \pm 444$, det$\times$accum
$573 \pm 202$, det$\times$prop $515 \pm 247$, prop$\times$accum $\mathbf{-330} \pm 242$.
Deux lectures et une seule est concluante : \textbf{il n'y a pas de paire dominante}, les deux
premières se tenant à $73$~M\euro{} pour des écarts-types de $708$ et $354$ ; mais les trois
paires porteuses passent toutes par la fréquence et font $86\,\%$ de l'ordre~2. Et
\textbf{prop$\times$accum est négative} sur les seize graines et les deux métriques : les deux
canaux de co-occurrence sont \emph{substituts}, un pilier déjà touché ne pouvant l'être deux
fois.

% =====================================================================================
\section{Le retour sur la conformité}
\src{\texttt{50}}

\paragraph{Le portage.} Détenir du capital coûte, \emph{chaque année}, un pourcentage du capital
immobilisé au titre de la marge de risque :
\begin{equation}
\text{portage} \;=\; \tau \times \Delta_{\text{DORA}}.
\end{equation}

\begin{res}
\textbf{Sortie} : à $\tau = 6\,\%$, $848$~M\euro/an ; au taux révisé $\tau = 4{,}75\,\%$,
$672$~M\euro/an. \textbf{L'écart entre les deux taux est publié plutôt que laissé au lecteur},
la chaîne restant au taux historique puisque la calibration est gelée.
\end{res}

\paragraph{Le bénéfice a deux composantes et une seule est monétisée.} La charge annuelle moyenne
passe de $3\,869$ à $622$~M\euro, soit une \textbf{sinistralité évitée de $3\,247$~M\euro/an},
qui vaut $3{,}8$~fois le portage. Ce qui est laissé de côté vaut donc près de quatre fois ce qui
est retenu.

\begin{cle}
\textbf{Le sens de la borne, et il était annoncé à l'envers.} Le retour se calcule comme un
capital divisé par un bénéfice annuel. Minorer le \emph{bénéfice} \textbf{majore} donc la
\emph{durée}. Le retour de $6$ à $35$ ans au taux historique, de $7$ à $45$ ans au taux révisé,
est un \textbf{MAJORANT} et non un plancher. Les deux composantes réunies donneraient $1$ à
$7$ ans, mais cette addition mêle un flux de compte de résultat à un coût du capital : c'est une
convention de retour sur investissement, pas une sortie de modèle.
\end{cle}

% =====================================================================================
\section{Robustesse : ce que chaque hypothèse coûte}

\subsection*{L'additivité des coûts au sein d'un sinistre}
\src{\texttt{74}}

\paragraph{Le relâchement.} La perte d'un sinistre touchant $k$ piliers vaut
\begin{equation}
L \;=\; \Bigl(\sum_{j \in \mathcal{C}} X_j\Bigr) \, k^{\,\theta - 1},
\qquad k = |\mathcal{C}|.
\end{equation}
La paramétrisation a deux propriétés qui la rendent lisible : $\theta = 1$ redonne
\emph{exactement} le modèle publié, ce qui sert de contrôle ; et $1^{\theta-1} = 1$, donc elle
laisse les sinistres mono-pilier là où ils sont et ne touche que la simultanéité.

\begin{center}\footnotesize
\renewcommand{\arraystretch}{1.15}
\begin{tabular}{@{}c r r r c@{}}
\toprule
$\boldsymbol{\theta}$ & \textbf{Conforme} & \textbf{Non conf.} &
$\boldsymbol{\Delta_{\text{DORA}}}$ & \textbf{/ publié} \\
\midrule
$0{,}70$ & $5\,435$ & $15\,141$ & $9\,706$ & $0{,}687$ \\
$1{,}00$ & $6\,049$ & $20\,188$ & $14\,139$ & $1{,}000$ \\
$1{,}30$ & $6\,972$ & $27\,459$ & $20\,487$ & $1{,}449$ \\
\midrule
\emph{élasticité} & $0{,}39$ & $0{,}95$ & $1{,}20$ & \\
\bottomrule
\end{tabular}
\end{center}

\begin{res}
\textbf{Ce que cela dit} : la plage vaut $76\,\%$ de l'écart publié et \textbf{$15$~fois} son
bruit de $\pm 734$~M\euro, donc l'effet est réel. \textbf{Et l'hypothèse n'est pas neutre entre
les deux états} : élasticités $0{,}95$ contre $0{,}39$, facteur $2{,}42$, parce que l'exposant ne
mord que sur les multi-piliers, majoritaires au seul état non conforme.\\
\textbf{Ce qui tient} : sur toute la plage l'écart garde son signe et son ordre de grandeur.
\textbf{Ce qui ne se déclare pas} : le \emph{sens}, mutualisation de la remédiation et saturation
de la capacité tirant en sens contraire. L'amplitude est \textbf{posée}, d'où la publication de
l'élasticité plutôt que de la plage.
\end{res}

\subsection*{La structure de la dépendance entre états}
\src{\texttt{69}}

\paragraph{Le modèle d'états.} Les états de conformité dérivent d'une latente à facteur commun
\begin{equation}
C^\star_j \;=\; \gamma \Theta \;+\; \sqrt{1 - \gamma^2}\, \varepsilon_j,
\qquad \gamma = 0{,}68,
\end{equation}
d'où une corrélation $\gamma^2 = 0{,}462$ entre deux piliers quelconques. \textbf{La dépendance
existe donc déjà et elle est forte} ; ce qui manque est sa \emph{structure}, échangeable au lieu
d'être propre à certains triplets. On ajoute un surcroît $\delta$ aux seules paires concernées,
par une matrice de corrélation dont la diagonale vaut un, ce qui préserve les marges par
construction.

\begin{res}
\textbf{Sortie} : $\delta$ admissible jusqu'à $0{,}38$ avant perte de positivité. Sur ce
balayage la loi des configurations bouge nettement, $P(\text{tous NC})$ de $0{,}068$ à $0{,}088$,
et \textbf{le capital espéré ne bouge que de $2$~M\euro, soit $0{,}02\,\%$}.\\
\textbf{Et la raison est structurelle, pas numérique} : le capital des $32$ configurations
s'ajuste par une forme \emph{additive} en indicateurs de pilier à $R^2 = 0{,}9945$, de
contributions $3\,524$ (P1), $2\,863$ (P4), $2\,022$ (P2), $1\,577$ (P3), $897$ (P5). L'espérance
d'une fonction additive ne dépend que des \emph{marges}, jamais de la dépendance : l'invariance
vaut donc pour \textbf{toute} structure à marges fixées, y compris non essayée.
\end{res}

\noindent\emph{Ce que cela ne dit pas} : l'invariance porte sur l'\textbf{espérance}. Un quantile
de la loi des configurations bougerait.

\subsection*{La mesure de queue, et son existence}
\src{\texttt{73}}

\paragraph{La CTE au niveau agrégé.} Elle est calculée en \emph{diagnostic}, jamais comme mesure
de couverture, le capital restant partout un $\mathrm{VaR}_{99,5\%}$ de la charge annuelle.

\begin{res}
\textbf{Sortie} : CTE de $5\,734$, $12\,607$ et $18\,183$~M\euro{} aux niveaux $95$, $99$ et
$99{,}5\,\%$. L'équivalence $\mathrm{CTE}_\beta = \mathrm{VaR}_{99,5\%}$ tombe à
$\beta = 97{,}73\,\%$, donc \textbf{la CTE à $95\,\%$ est moins prudente} que la VaR
réglementaire, d'un facteur $1{,}46$. Rapport $\mathrm{CTE}/\mathrm{VaR} = 2{,}17$ contre
l'asymptote $1/(1-\xi) = 2{,}47$.
\end{res}

\begin{cle}
\textbf{L'argument décisif est l'existence, pas la proportion.} Sous la calibration alternative
$\hat\xi = 1{,}033 > 1$, l'espérance de la sévérité est infinie et \textbf{la mesure n'existe
pas}. Une couverture qui cesse d'être définie selon la source de sévérité ne peut pas porter un
capital. À utiliser plutôt que « la CTE est excessive », qui est un argument de degré.
\end{cle}

\subsection*{Le niveau de l'intervalle reporté}
\src{\texttt{72}}

\paragraph{Le passage de $90$ à $95\,\%$ ne rejoue aucun bootstrap.} Il applique le rapport des
quantiles normaux aux deux demi-largeurs \emph{séparément}, ce qui préserve l'asymétrie :
\begin{equation}
[\,b_{95},\,h_{95}\,] \;=\;
\Bigl[\, p - \tfrac{z_{95}}{z_{90}}(p - b_{90}),\;\; p + \tfrac{z_{95}}{z_{90}}(h_{90} - p) \,\Bigr],
\qquad \frac{z_{95}}{z_{90}} = \frac{1{,}96}{1{,}645} = 1{,}1915.
\end{equation}

\begin{res}
\textbf{Sortie} : $\xi$ passe à $[0{,}2487\,;0{,}8765]$, soit un facteur $3{,}52$ ; le capital à
$[2\,727\,;34\,942]$~M\euro, facteur $12{,}81$.\\
\textbf{Deux constats} : le changement de niveau ne déplace \textbf{aucune calibration}, ni le
point ni les paramètres, seulement les bornes publiées ; et il \textbf{ne corrige pas} la
sous-couverture, qui reste un déficit constant de $-3{,}2$~points ($86{,}8 \pm 0{,}8\,\%$ réels
pour $90$ nominaux, $91{,}8 \pm 0{,}6\,\%$ pour $95$).
\end{res}

\subsection*{La transposition à l'échelle d'une entité}
\src{\texttt{60}, \texttt{70}}

\paragraph{Ce que la renormalisation proportionnelle suppose.} Renormaliser au prorata du capital
de marché affirme une \textbf{élasticité un} à la taille. Les deux canaux estimés en sont loin :
fréquence $+0{,}0744$ (écart-type $0{,}0098$, $z = 7{,}59$), sévérité $0{,}087$ d'intervalle
$[0{,}026\,;0{,}148]$.

\begin{res}
\textbf{Sortie} : l'élasticité globale mesurée vaut $0{,}204$ contre $1{,}000$ supposée. Une
petite entité reçoit $0{,}633$ contre $0{,}100$ sous proportionnalité, une grande $1{,}621$ contre
$10{,}00$.\\
\textbf{Les deux approches se trompent en sens opposés} : la proportionnalité sous-estime une
petite entité puisqu'elle rétrécit une sévérité qui ne rétrécit pas, le modèle la majore puisqu'il
ne la plafonne pas. Substituer l'une à l'autre changerait le signe du défaut sans le déclarer.
\end{res}

\subsection*{Les sensibilités, et la fragilité n'est pas où on l'attend}
\src{\texttt{15}, \texttt{30}, \texttt{43}}

\begin{res}
\textbf{Tornado} : le gain de propagation $g$ de $0{,}5$ à $1{,}0$ déplace le surcoût de
$5\,517$ à $6\,579$~M\euro, quand le \emph{seuil} de la loi de queue le déplace de $4\,862$ à
$12\,944$ et la fréquence de $4\,056$ à $8\,684$. \textbf{La propagation est le levier le moins
sensible des quatre} : la fragilité est dans l'ajustement de valeurs extrêmes, pas dans la
contagion.\\
\emph{Garde} : ces niveaux sont ceux de la base du tornado ($6\,640$~M\euro), qui n'est pas
l'écart à quatre canaux de $14\,139$.
\end{res}

\begin{res}
\textbf{Coût de l'ignorance directionnelle} : $336$~M\euro{} de largeur à $t = 0{,}25$, $695$ à
$0{,}50$, $1\,083$ à $0{,}75$, $1\,839$ à $t = 1$. \textbf{L'ignorance se paye linéairement et son
coût maximal est connu d'avance}, ce qui est plus favorable qu'un intervalle de confiance
ordinaire. Le point d'expert ($8\,110$, soit $+53{,}7\,\%$ du socle) est un point de l'ensemble et
non sa mesure, et la direction nulle donne déjà $+51{,}5\,\%$ : l'essentiel du surcoût vient de la
\emph{co-occurrence}, qui est identifiée, non de la direction.
\end{res}

\subsection*{Ce que les séquences ordonnées permettent, et ce qu'elles interdisent}
\src{\texttt{75}}

\begin{res}
\textbf{Matière disponible} : $22$ transitions codées sur $10$ incidents, réparties sur $6$ paires
distinctes, d'où $9$ chemins de longueur deux et $4$ séquences distinctes.\\
\textbf{Pourquoi la question d'asymétrie P1\,/\,P2 ne se tranche pas par les triplets} : P2
n'émet \textbf{jamais}, $0$ fois sur $22$, et P1 n'est jamais atteint. Le support est vide.\\
\textbf{Pourquoi le test n'a aucune puissance} : il exige d'un même pilier qu'il soit atteint
\emph{et} qu'il ait deux successeurs distincts ; aucun ne remplit les deux, et les deux manques
sont exclusifs. Il faudrait de $158$ à $589$ chemins par un même pilier qualifié.\\
\textbf{Acquis} : la chaîne \textbf{n'est pas sans mémoire}, l'auto-évitement gonflant la loi du
successeur d'un facteur $1{,}211$ en moyenne, et c'est une borne basse.\\
\textbf{Précaution} : le graphe agrégé est acyclique, mais $42$ des $64$ orientations le seraient
aussi, soit $65{,}6\,\%$. L'acyclicité n'est donc \textbf{pas} un résultat et ne renforce pas la
frontière d'identification.
\end{res}

% =====================================================================================
\section{Les défauts connus, chiffrés plutôt que corrigés}
\src{\texttt{08b}, \texttt{41}, \texttt{47}, \texttt{57}, \texttt{74}}

\paragraph{Pourquoi ils ne sont pas corrigés.} La calibration est gelée : une correction
déplacerait des centaines de nombres publiés sans qu'aucun déplacement soit attribuable à la
correction, et la piste d'audit serait perdue pour un gain immatériel. Une limite chiffrée,
signée et recalculée à chaque exécution vaut mieux qu'un nombre corrigé en silence.

\begin{center}\footnotesize
\renewcommand{\arraystretch}{1.2}
\begin{tabular}{@{}p{5.0cm} p{2.4cm} p{8.4cm}@{}}
\toprule
\textbf{Écart} & \textbf{Sens} & \textbf{Taille} \\
\midrule
\multicolumn{3}{@{}l}{\emph{Défauts d'estimation, tous deux anti-conservateurs}} \\
Taux de dépassement gelé $p_u$ & sous-estime & $+2{,}4\,\%$ de quantile manquant, soit $15{,}8$~M\euro, ou $2{,}5\,\%$ de la largeur de l'intervalle de la même VaR \\
Couverture réelle de l'intervalle & sous-estime & $86{,}8\,\%$ au lieu de $90$ ; demi-largeur à élargir d'environ $10\,\%$ \\
\midrule
\multicolumn{3}{@{}l}{\emph{Choix posés, tous prudents}} \\
Queue de sévérité $\xi$ & majore & $0{,}90$ posé contre $0{,}5954$ estimé \\
Charge systémique & majore & $0{,}60$ posé, au-delà de la borne haute empirique $0{,}223$ \\
Sur-dispersion $\varphi$ & majore & $9{,}20$ posé contre $2{,}24$ mesuré \\
\midrule
\multicolumn{3}{@{}l}{\emph{Hypothèse dont le sens n'est pas déterminable}} \\
Additivité des coûts multi-piliers & \textbf{indéterminé} & élasticité $1{,}20$ ; $[9\,706\,;20\,487]$~M\euro{} sur l'amplitude posée \\
\bottomrule
\end{tabular}
\end{center}

\vspace{3pt}
\begin{cle}
\textbf{La lecture qui compte.} Les deux seuls écarts \emph{involontaires} sous-estiment le
capital, de quelques pour cent chacun ; les trois choix \emph{délibérés} le majorent, et de
beaucoup plus. \textbf{Le mémoire ne se protège donc pas derrière ses erreurs} : sa prudence vient
de choix assumés qu'il nomme comme tels, et corriger les deux premiers ne changerait aucune
conclusion. La dernière ligne est la seule dont le \emph{sens} ne se déclare pas.
\end{cle}

\end{document}
